% 原书印刷页85--86；本版PDF页94--95。
% 注意底本在第213题后将第214题误排为“224”，后文又继续215。
\thermophaseproblem{205}{85--86}
标准大气压下，
把水过冷到 $t=-20^\circ\mathrm C$。
随后体系自发转变到相平衡态时，
有多大质量分数 $\eta$ 的水结冰？
若希望过冷水全部结冰，
需要初始过冷至多少温度？
\exsolution
设液态过冷水质量为 $m$，
初温 $t<0^\circ\mathrm C$。
按绝热转变模型，结冰放出潜热
$\eta m q_m$，使尚未冻结的物质
从 $t$ 升至平衡凝固点 $0^\circ\mathrm C$。
以液态水比热 $c\simeq4.18\,\mathrm{kJ/(kg\,K)}$
近似，
热平衡为
\[
\eta m q_m\simeq mc|t|,
\]
故
\[
\boxed{\eta\simeq\frac{c|t|}{q_m}}.
\]
取冰融化潜热
$q_m\simeq334\,\mathrm{kJ/kg}$、
$|t|=20\,\mathrm K$：
\[
\boxed{\eta\simeq\frac{4.18\times20}{334}
\simeq0.25=25\%}.
\]
若过冷足以使全部质量结冰，
临界条件 $\eta=1$ 给
\[
\boxed{t_{\rm full}\simeq-\frac{q_m}{c}
\simeq-80^\circ\mathrm C}.
\]
\excheck $-80^\circ\mathrm C$ 是常比热、
无外部换热模型的\emph{形式阈值}，
不声称现实纯水在常压下
能稳定维持到该过冷程度。

\thermophaseproblem{206}{86}
在 $0^\circ\mathrm C$ 附近，
压强升高 $\Delta p=1.00\,\mathrm{atm}$。
冰的质量比体积比液态水大
$\Delta v=0.091\,\mathrm{cm^3/g}$。
求冰的熔点变化。
\exsolution
沿固液共存线，
Clapeyron 方程
\[
\frac{\mathrm dp}{\mathrm dT}
=\frac{q_m}{T(v_{\rm water}-v_{\rm ice})}
=-\frac{q_m}{T\Delta v}.
\]
在压强变化较小时
\[
\boxed{\Delta T\simeq
-\frac{T\Delta v}{q_m}\Delta p}.
\]
将 $\Delta v=0.091\,\mathrm{cm^3/g}
=9.1\times10^{-5}\,\mathrm{m^3/kg}$、
$T\simeq273.15\,\mathrm K$、
$q_m\simeq3.34\times10^5\,\mathrm{J/kg}$ 代入，
\[
\boxed{\Delta T\simeq-7.5\times10^{-3}\,\mathrm K
=-7.5\,\mathrm{mK}}.
\]
\excheck 因液态水比冰密，压强升高利于液态，
熔点降低。本版纯文本提取将体积单位
$\mathrm{cm^3/g}$ 误作 $\mathrm{cm^2/g}$，
已按物理量纲澄清。

\thermophaseproblem{207}{86}
在标准气压下，
如果压强降低 $\Delta p=3.2\,\mathrm{kPa}$
会使水的沸点降低 $\Delta T=0.9\,\mathrm K$，
求原来饱和水蒸汽的质量比体积。
\exsolution
Clapeyron 方程
$\mathrm dp/\mathrm dT=
q_v/[T(v_v-v_\ell)]$。
在沸点附近
$v_v\gg v_\ell$，
故
$\Delta p/\Delta T
\simeq q_v/(T v_v)$。
取压强、温度变化的模：
\[
\boxed{v_v\simeq\frac{q_v\Delta T}{T\Delta p}}.
\]
用
$q_v\simeq2.26\times10^6\,\mathrm{J/kg}$、
$T\simeq373\,\mathrm K$：
\[
\boxed{v_v\simeq
\frac{2.26\times10^6(0.9)}{373(3.2\times10^3)}
\simeq1.7\,\mathrm{m^3/kg}}.
\]
\excheck 这是单位\emph{质量}蒸汽的体积，
而非一摩尔蒸汽的体积；
结果与 $100^\circ\mathrm C$ 附近饱和蒸汽量级一致。

\thermophaseproblem{208}{86}
假设饱和水蒸气为理想气体，
求温度
$101.1^\circ\mathrm C$ 时的饱和蒸气压。
\exsolution
以 $T_0=373.15\,\mathrm K$
处 $p_0=1\,\mathrm{atm}$ 为参考，
温度增量
$\Delta T=1.1\,\mathrm K$。
在忽略液态比体积和汽化潜热温度依赖的近似下，
\[
\frac{\mathrm d\ln p}{\mathrm dT}
=\frac{q_vM}{RT^2}.
\]
对小 $\Delta T$，用一阶展开：
\[
p(T_0+\Delta T)
\simeq p_0\left[
1+\frac{q_vM}{RT_0^2}\Delta T\right].
\]
代入
$q_v\simeq2.26\times10^6\,\mathrm{J/kg}$、
$M=0.018\,\mathrm{kg/mol}$：
\[
\boxed{p\simeq1.04\,\mathrm{atm}}.
\]
\excheck 本题温差仅1.1 K，
可以用线性近似；
若温差较大，则应对 $T^{-2}$ 积分。

\thermophaseproblem{209}{86}
在 $100^\circ\mathrm C$ 时，
少量水与其饱和蒸汽共存于密闭容器。
温度升高 $\Delta T=1.5\,\mathrm K$，
在水的液相比体积可以忽略、
饱和蒸汽为理想气体的条件下，
求饱和蒸汽质量增加的百分比。
\exsolution
容器总体积固定，液相占容积很小。
饱和蒸汽质量为
$m_v=p_{\rm sat}VM/(RT)$，所以
\[
\frac{\mathrm dm_v}{m_v}
=\frac{\mathrm dp_{\rm sat}}{p_{\rm sat}}
-\frac{\mathrm dT}{T}.
\]
Clausius--Clapeyron 近似给
$\mathrm d\ln p_{\rm sat}/\mathrm dT
=q_vM/(RT^2)$。
因此
\[
\frac{\Delta m_v}{m_v}
\simeq
\left(\frac{q_vM}{RT}-1\right)\frac{\Delta T}{T}.
\]
代入 $T=373.15\,\mathrm K$、
$\Delta T=1.5\,\mathrm K$：
\[
\boxed{\frac{\Delta m_v}{m_v}\simeq0.05
\approx5\%}.
\]
\excheck 必须同时计入
饱和蒸汽压随温度上升和
理想气体密度中 $1/T$ 的减小因子。

\thermophaseproblem{210}{86}
已知饱和蒸汽压强在温度 $T_0$ 时为 $p_0$。
假设汽化比潜热 $q_v$ 与温度无关、
液相比体积远小于蒸汽比体积、
饱和蒸汽为理想气体，
求饱和蒸汽压强函数 $p(T)$，
并说明这些近似的适用范围。
\exsolution
Clapeyron 方程
$\mathrm dp/\mathrm dT
=q_v/[T(v_v-v_\ell)]$。
取 $v_v\gg v_\ell$ 且
$v_v=RT/(Mp)$：
\[
\frac1p\frac{\mathrm dp}{\mathrm dT}
=\frac{q_vM}{RT^2}.
\]
在 $T_0$ 到 $T$ 上积分：
\[
\ln\frac{p(T)}{p_0}
=\frac{q_vM}{R}
\int_{T_0}^{T}\frac{\mathrm dT'}{T'^2}
=\frac{q_vM}{R}
\left(\frac1{T_0}-\frac1T\right).
\]
故
\[
\boxed{p(T)=p_0
\exp\left[\frac{q_vM}{R}
\left(\frac1{T_0}-\frac1T\right)\right]}.
\]
\excheck 近似适用于远低于临界温度的
\emph{较窄温区}；
靠近临界点时蒸汽与液态比体积接近，
潜热随温度显著变动，上式失效。

\thermophaseproblem{211}{86}
处于初始标准状况的冰，
被压缩到 $p=640\,\mathrm{atm}$。
假设冰的熔点下降与压强升高成正比，
水与冰的质量比体积差
$\Delta v=0.09\,\mathrm{cm^3/g}$
（液水比冰小）。
求有多少质量分数的冰熔化。
\exsolution
由于压强升高降低熔点，
在绝热压缩与相平衡近似下，
冰体温度需随新的融化温度降低，
一部分冰熔化并吸收潜热。
由 Clapeyron 近似
\[
|\Delta T|\simeq
\frac{T\Delta v}{q_m}\,\Delta p.
\]
降温释放的冰显热约为
$mc_{\rm ice}|\Delta T|$，
用于熔化质量 $\eta m$ 的潜热：
$\eta m q_m\simeq mc_{\rm ice}|\Delta T|$。
因此
\[
\boxed{\eta\simeq
\frac{c_{\rm ice}T\Delta v\,\Delta p}{q_m^2}}.
\]
取
$c_{\rm ice}\simeq2.1\times10^3\,\mathrm{J/(kg\,K)}$、
$T\simeq273\,\mathrm K$、
$\Delta p\simeq639\,\mathrm{atm}$、
$q_m\simeq3.34\times10^5\,\mathrm{J/kg}$，
\[
\boxed{\eta\simeq0.03=3\%}.
\]
\excheck 这里使用了\emph{压缩近似绝热}
与熔点线性下降假设；
若大量压缩功以热形式内化，需进一步修正能量平衡。

\thermophaseproblem{212}{86}
二氧化碳三相点附近的饱和蒸汽压
满足以10为底的经验关系
$\log_{10}p=a-b/T$，
其中 $p$ 以 atm 为单位。
升华支路：
$a_s=9.05,\ b_s=1.80\,\mathrm{kK}$；
汽化支路：
$a_v=6.78,\ b_v=1.31\,\mathrm{kK}$。
求：(a) 三相点温度和压强；
(b) 三相点附近的升华、汽化、熔化质量潜热。
\exsolution
三相点两支路压强相等：
\[
a_s-\frac{b_s}{T_t}
=a_v-\frac{b_v}{T_t}.
\]
解得
\[
\boxed{T_t=\frac{b_s-b_v}{a_s-a_v}
=\frac{1800-1310}{9.05-6.78}
\simeq216\,\mathrm K}.
\]
任一支路回代：
\[
\boxed{p_t=10^{\,a_s-b_s/T_t}
\simeq5.1\,\mathrm{atm}}.
\]
由理想蒸汽和 Clausius--Clapeyron 近似
$\mathrm d\ln p/\mathrm dT
=L_{\rm mol}/(RT^2)$。
题给为\emph{十进制对数}：
\[
\frac{\mathrm d\ln p}{\mathrm dT}
=(\ln10)\frac b{T^2},
\quad L_{\rm mol}=R(\ln10)b.
\]
故每单位质量潜热
$q=R(\ln10)b/M$，
二氧化碳 $M\simeq0.044\,\mathrm{kg/mol}$：
\[
\boxed{q_s\simeq0.78\,\mathrm{kJ/g}},\qquad
\boxed{q_v\simeq0.57\,\mathrm{kJ/g}}.
\]
由三相点焓变相加，
$q_s=q_m+q_v$，所以
\[
\boxed{q_m=q_s-q_v\simeq0.21\,\mathrm{kJ/g}}.
\]
\excheck 必须把 $\log_{10}$ 转为自然对数，
带上 $\ln10$，否则潜热会被低估约2.3倍。

\thermophaseproblem{213}{86}
质量 $m=1.00\,\mathrm{kg}$ 的水，
由 $t_1=10^\circ\mathrm C$ 加热至
$t_2=100^\circ\mathrm C$，
再于沸点全部汽化。
求系统的熵增。
\exsolution
先将液态水由
$T_1=283.15\,\mathrm K$
加热至
$T_2=373.15\,\mathrm K$，
视比热 $c$ 为常数，
熵增为
$m c\ln(T_2/T_1)$。
随后在温度 $T_2$ 等温汽化，
熵增为
$m q_v/T_2$。
两段相加：
\[
\boxed{\Delta S=
m\left[c\ln\frac{T_2}{T_1}+\frac{q_v}{T_2}\right]}.
\]
取水 $c\simeq4.18\,\mathrm{kJ/(kg\,K)}$、
$q_v\simeq2.26\,\mathrm{MJ/kg}$，
\[
\boxed{\Delta S\simeq7.2\,\mathrm{kJ/K}}.
\]
\excheck 绝对温度必须用 K；
先升温后相变的两段熵贡献不能遗漏其中一项。

\thermophaseproblem{214}{86}
冰初温为 $t_1=0^\circ\mathrm C$，
先全部融化，再把水加热至
$t_2=100^\circ\mathrm C$，最后完全汽化。
求系统\emph{单位质量}的熵增。
\exsolution
熔化发生在
$T_m\simeq273.15\,\mathrm K$，
其单位质量熵增
$q_m/T_m$。
液态水升温至
$T_b\simeq373.15\,\mathrm K$，
熵增
$c\ln(T_b/T_m)$。
在沸点汽化，
熵增
$q_v/T_b$。
总单位质量熵增：
\[
\boxed{\Delta s=
\frac{q_m}{T_m}
+c\ln\frac{T_b}{T_m}
+\frac{q_v}{T_b}}.
\]
用
$q_m\simeq334\,\mathrm{J/g}$、
$c\simeq4.18\,\mathrm{J/(g\,K)}$、
$q_v\simeq2257\,\mathrm{J/g}$，
\[
\boxed{\Delta s\simeq8.6\,\mathrm{J/(g\,K)}}.
\]
\excheck \textbf{原书题号排印：}
在本版正文 PDF 印刷页86，
第213题之后此题被误印为“224”，
下题紧接着却是“215”；书后答案亦对应第214题。
为使题库和第213、215题连续，
本稿按\emph{第214题}录入并保留此处勘误。
